\documentclass[11pt,compress,t,notes=noshow, aspectratio=169, xcolor=table]{beamer}

\usepackage{../../style/lmu-lecture}
% Defines macros and environments
\input{../../style/common.tex}

\title{Applied Machine Learning}
% \author{LMU}
%\institute{\href{https://compstat-lmu.github.io/lecture_iml/}{compstat-lmu.github.io/lecture\_iml}}
\date{}

\begin{document}
\newcommand{\titlefigure}{slides/03-ensembles-stacking/figure/voting.png}
\newcommand{\learninggoals}{
\item Why ensembling is useful
\item Different techniques
\item Stacking}

\lecturechapter{Ensembling}
\lecture{Applied Machine Learning}

% ------------------------------------------------------------------------------

\begin{vbframe}{Ensembling in Machine Learning}
\begin{itemize}
\item Techniques that combine predictions of multiple base models
\item Goal: Improve predictive performance and model robustness
\item Idea: Errors of a single model will likely be compensated by other models, and as a result, the overall prediction performance of the ensemble would be better than that of a single model
\end{itemize}
\vspace{1em}
\begin{columns}
    \begin{column}{0.5\textwidth}
        \begin{figure}
            \centering
            \includegraphics{slides/03-ensembles-stacking/figure/ensemble_smiley.png}
            %\caption{Caption}
            % from https://wires.onlinelibrary.wiley.com/doi/full/10.1002/widm.1249?saml_referrer
            \label{fig:ensemble_smiley}
        \end{figure}
    \end{column}
    \begin{column}{0.5\textwidth}
        \begin{figure}
            \centering
            \includegraphics{slides/03-ensembles-stacking/figure/ensembles_hs.png}
            %\caption{Caption}
            %from https://web.engr.oregonstate.edu/~tgd/publications/mcs-ensembles.pdf
            \label{fig:ensembles_hs}
        \end{figure}
    \end{column}
\end{columns}
\end{vbframe}


% ------------------------------------------------------------------------------

\begin{vbframe}{Ensemble Methods}
Different approaches:
\begin{itemize}
\item Model Averaging
\item Bagging
%\item Boosting
\item Stacking
\end{itemize}
\vspace{1em}
Higher level mechanisms:
\begin{itemize}
    \begin{itemize}
        \item during base model training (how to generate different models?)
        \begin{itemize}
            \item data manipulation (e.g., bagging)
            \item partitioning of data
            \item learning algorithm manipulation (e.g., randomly select splits)
            \item output manipulation (error-correcting output codes)
            \item hybridization of the above
        \end{itemize}
        \item during prediction (how to combine the predictions of different models?)
        \begin{itemize}
            \item output manipulation
            \item weighting methods
            \item meta-learning
        \end{itemize}
        \item homogeneous vs. heterogeneous
        \begin{itemize}
            \item same types of models vs. different models
        \end{itemize}
    \end{itemize}
\end{itemize}
 \end{vbframe}

% ------------------------------------------------------------------------------

\begin{vbframe}{Notation}
\begin{itemize}
\item $\mathcal{D}$ training data
\item $\{f_{1}, \ldots, f_{m}\} \subseteq \mathcal{H}$ set of $m$ models trained on $\mathcal{D}$
\item $f^{[M]} = G(f_{1}, \ldots, f_{m})$ ensemble model obtained via an aggregating function $G: \mathcal{H}^m \rightarrow \mathcal{H}$
\end{itemize}
\begin{figure}
    \centering
    \includegraphics[width=0.7\textwidth]{slides/03-ensembles-stacking/figure/rf_majvot_averaging.png}
    %\caption{Caption}
    % from i2ml
    \label{fig:rf_mjavot_averaging}
\end{figure}
\end{vbframe}

% ------------------------------------------------------------------------------

\begin{vbframe}{Model Averaging}

Definition: Combine predictions of multiple (different) models to make a final prediction\\
\begin{columns}
    \begin{column}{0.5\textwidth}
        \vspace{-1em}
        \begin{figure}
            \centering
            \includegraphics{slides/03-ensembles-stacking/figure/voting.png}
            %\caption{Caption}
            % from https://sebastianraschka.com/pdf/lecture-notes/stat479fs19/07-ensembles__notes.pdf
            \label{fig:voting}
        \end{figure}
    \end{column}
    \begin{column}{0.5\textwidth}
        Steps:
        \begin{enumerate}
            \item Base model training
            \item Combine predictions
        \end{enumerate}
    \end{column}
\end{columns}
\vspace{1em}
Combine predictions: majority voting (hard) / model averaging (soft)
\begin{itemize}
    \item $f^{[M]}(\mathbf{x}) = \text{mode}(\{f_{1}(\mathbf{x}), \ldots, f_{m}(\mathbf{x})\})$ assuming hard label predictions    
    \item $f^{[M]}(\mathbf{x}) = \frac{1}{M} \sum_{i = 1}^{M} f_{i}(\mathbf{x})$ assuming probability predictions
\end{itemize}
\end{vbframe}

% ------------------------------------------------------------------------------

\begin{vbframe}{Weighted Model Averaging}

We can also introduce weights $\mathbf{w} = (w_{i}, \ldots, w_{m}), \text{usually: } 0 \le w_{i} \le 1, \sum_{i=1}^{M} w_{i} = 1$ accounting for the contribution of each individual model to the prediction of the ensemble: $f_{\mathbf{w}}^{[M]}(\mathbf{x}) = \frac{1}{M} \sum_{i = 1}^{M} w_{i} f_{i}(\mathbf{x})$\\
\vspace{1em}
This gives rise to an optimization problem to find the optimal weight vector so that the generalization error of the ensemble is optimal:
$\argmin_{\mathbf{w} \in \mathcal{W}} \hat{\text{GE}}(f_{\mathbf{w}}^{[M]})$\\
\vspace{1em}
Can be optimized via:
\begin{itemize}
    \item linear model
    \item greedy ensemble selection
    \item ...
\end{itemize}

\end{vbframe}

% ------------------------------------------------------------------------------

\begin{vbframe}{Bagging}

Definition: Bagging (stands for Bootstrap Aggregating), training multiple instances of the same base model on different subsets of the training data, obtained through random sampling with replacement.\\
\begin{columns}
    \begin{column}{0.5\textwidth}
        \vspace{-1em}
        \begin{figure}
            \centering
            \includegraphics{slides/03-ensembles-stacking/figure/bagging.png}
            %\caption{Caption}
            % from https://sebastianraschka.com/pdf/lecture-notes/stat479fs19/07-ensembles__notes.pdf
            \label{fig:baggiing}
        \end{figure}
    \end{column}
    \begin{column}{0.5\textwidth}
        Steps:
        \begin{enumerate}
            \item Random sampling with replacement to generate different train data sets
            \item Base model training
            \item Combine predictions
        \end{enumerate}
    \end{column}
\end{columns}
\vspace{1em}
Popular use case: Construction of a random forest using decision trees as base models

\end{vbframe}

% ------------------------------------------------------------------------------

% ------------------------------------------------------------------------------

\begin{vbframe}{Bagging}

Why subsets of training data obtained through random sampling with replacement (input manipulation):
\begin{itemize}
    \item Bias vs. variance trade-off
    \begin{itemize}
        \item By training each model on slightly different data the variance of the predictions are reduced
        \item The training data is reused many times so bias can be slightly higher  
    \end{itemize}
    \item Effective for cases where small changes in the training set may result in a completely different model (instable learners)
\end{itemize}
\vspace{1em}
Note that a random forest does not only employ input manipulation via bagging but also learning algorithm manipulation (sampling splits).

\end{vbframe}

% ------------------------------------------------------------------------------

\begin{vbframe}{Boosting}
ToDo: Don't cover boosting here detailed like the other topics, because boosting is more or less a learning method (similarly as a random forest which we do not cover except for the example in bagging).
Briefly mention this and examples adaboost, xgboost, ... and refer to other lectures
%Definition: Boosting is an iterative technique that adjusts the weights of the training instances based on their misclassification rate or residual errors, and trains subsequent base models to focus on the misclassified or poorly predicted instances.

%Steps:

%begin{itemize}
%\item Initial model training:
%\item Weight update
%\item Base model training using updated weights
%\item Repeat weight update
%\item ...
%\item Combine predictions
%\end{itemize}
\end{vbframe}

% ------------------------------------------------------------------------------

\begin{vbframe}{Stacking}
Definition: Stacking (also stacked generalization or meta-ensemble), is an ensemble technique that involves training multiple base models on the same data and using their predictions as input features for a meta-model that makes the final ensemble prediction.

\begin{columns}
    \begin{column}{0.5\textwidth}
        \vspace{-1em}
        \begin{figure}
            \centering
            \includegraphics[width=0.9\textwidth]{slides/03-ensembles-stacking/figure/stacking.png}
            %\caption{Caption}
            % from https://sebastianraschka.com/pdf/lecture-notes/stat479fs19/07-ensembles__notes.pdf
            \label{fig:baggiing}
        \end{figure}
    \end{column}
    \begin{column}{0.5\textwidth}
        Steps:
        \begin{enumerate}
            \item Base model training
            \item Collect predictions
            \item Meta model training
        \end{enumerate}
    \end{column}
\end{columns}
\vspace{1em}

In contrast to the other methods, stacking already combines model predictions during training.

\end{vbframe}

% ------------------------------------------------------------------------------

\begin{vbframe}{Stacking Cross-validated}

\begin{columns}
    \begin{column}{0.5\textwidth}
        \vspace{-1em}
        \begin{figure}
            \centering
            \includegraphics[width=0.9\textwidth]{slides/03-ensembles-stacking/figure/stacking_cv.png}
            %\caption{Caption}
            % from https://sebastianraschka.com/pdf/lecture-notes/stat479fs19/07-ensembles__notes.pdf
            \label{fig:baggiing}
        \end{figure}
    \end{column}
    \begin{column}{0.5\textwidth}
        Using cross-validated predictions of the base models
        \begin{itemize}
            \item improve generalization
            \item reduce overfitting
            \item enhances base model diversity
            \item minimizes information leakage
        \end{itemize}
    \end{column}
\end{columns}
\end{vbframe}

% ------------------------------------------------------------------------------

\begin{vbframe}{Multilevel Stacking}
So far: train L0 base models, use predictions to train the L1 meta model.\\
Now: train multiple meta models at L1, use predictions to train final L2 meta model 

\begin{columns}
    \begin{column}{0.5\textwidth}
        \vspace{-1em}
        \begin{figure}
            \centering
            \includegraphics[width=0.9\textwidth]{slides/03-ensembles-stacking/figure/stacking.png}
            %\caption{Caption}
            % from https://sebastianraschka.com/pdf/lecture-notes/stat479fs19/07-ensembles__notes.pdf
            \label{fig:baggiing}
        \end{figure}
    \end{column}
    \begin{column}{0.5\textwidth}
        Steps:
        \begin{enumerate}
            \item L0 base model training
            \item Collect predictions of L0 base models
            \item L1 meta model training 
            \item Collect predictions of L1 models
            \item L2 final meta model training
        \end{enumerate}
    \end{column}
\end{columns}
\vspace{1em}
\end{vbframe}

% ------------------------------------------------------------------------------

\begin{vbframe}{Meta Models}
Popular choices of meta models in stacking include:
\begin{itemize}
    \item linear / logistic regression 
    \item random forest (especially in the case of multilevel stacking)
\end{itemize}
\end{vbframe}

% ------------------------------------------------------------------------------

\begin{vbframe}{Open topics}
\begin{itemize}
    \item calibration of ensembles
    \item how detailed bias variance recap
    \item diversity in ensembling
    \item HPO
\end{itemize}
\end{vbframe}

% ------------------------------------------------------------------------------


\endlecture
\end{document}
